% =============================================================================
%  q_ramification_along_VF.tex
%
%  Research suggestions:  controlling the ramification of the symmetrisation
%  map  q : (cP_{G/H})^{(d)} -> (cP_{G/H})^{[d]}  (equivalently the field
%  extension  K_{C'} / F )  ALONG the place  V_F = V_{C'}|_F ,  and using the
%  global hypotheses
%
%        ram(cP -> U) bounded by d        and        G = Z/p^r ,  r < d
%
%  to supply the "missing global input" that the inductive proof of
%  \Cref{theorem:ramification_after_blowup_is_tame_for_p_power} lacks
%  (the red TODO at RamificationAfterBlowupIsTame.tex:671, and the conclusion
%  of a_wrong_path.tex).
%
%  Compile:  latexmk -pdf q_ramification_along_VF.tex   (uses biber)
% =============================================================================
\documentclass[11pt]{article}
\usepackage{geometry}\geometry{margin=1in}
\usepackage[T1]{fontenc}
\usepackage{ifthen}
\usepackage{booktabs}

\setlength{\parindent}{1em}
\setlength{\parskip}{0.5em}

\input{includes}
\input{MacrosAndFigures}

% --- local fallbacks for symbols not in the thesis macro file ---------------
\providecommand{\PP}{\mathbb{P}}
\providecommand{\fm}{\mathfrak{m}}
\providecommand{\Bl}{\mathrm{Bl}}

\usepackage{csquotes}
\usepackage[
backend=biber,
style=alphabetic,
sorting=ynt
]{biblatex}
\addbibresource{Bib.bib}

\newboolean{ThesisIntro}
\setboolean{ThesisIntro}{false}

\title{Controlling the ramification of $q$ along $V_F$\\[2pt]
       {\large how the global bound $\mathrm{ram}(\cP)\le d$ together with
        $G=\Z/p^r$, $r<d$, can supply the input the induction is missing}}
\author{Assaf Marzan}
\date{\today}

\begin{document}
\maketitle

\noindent\textbf{Purpose.}
The inductive proof of
\Cref{theorem:ramification_after_blowup_is_tame_for_p_power}
($\cP^{[\deg\ndls]}$ unramified at $\eta_\ndls$ for $G=\Z/p^r$) stalls at one
point (the red note ``\emph{not clear how to finish without analysing the
ramification even further}'' in \texttt{RamificationAfterBlowupIsTame.tex}).
The post-mortem in \texttt{a\_wrong\_path.tex} localised the stall to a single,
sharp phenomenon: after reducing to
\[
\textbf{Goal:}\qquad E/F \text{ is unramified at } V_F,\qquad I(v_E\mid V_F)=1,
\]
the hypothesis (H2) only gives $I(v_\star\mid V_{C'})=1$, and the two are linked
by an injection
\begin{equation}\label{eq:inj}
I(v_\star\mid V_{C'})\;\hookrightarrow\;I(v_E\mid V_F)
\end{equation}
that need not be onto. The cokernel is \emph{exactly} the ramification of
\[
K_{C'}/F \quad\text{at}\quad V_{C'}\mid V_F,
\]
i.e.\ the ramification that the symmetrisation map
$q:(\cP_{G/H})^{(d)}\to(\cP_{G/H})^{[d]}$ introduces at the point at infinity.
When that ramification is wild it can \emph{absorb} the ramification of $E/F$,
and \texttt{a\_wrong\_path.tex} exhibits a local coincidence
$\widehat{K_{C'}}=\widehat E$ making (H1)--(H3) hold while the Goal fails. Its
closing remark names the cure:

\begin{quote}\itshape
``ruling out such coincidences for a given $(\cP,d)$ is exactly the global input
the inductive argument lacks.''
\end{quote}

This note collects concrete suggestions for supplying that global input by
\emph{measuring the ramification of $q$ along $V_F$} and bounding it by the two
global hypotheses we have not yet used:
\begin{equation}\label{eq:hyp}
\textbf{(B)}\ \ \mathrm{ram}(\cP\to U)\ \text{bounded by}\ d\ \ (\text{i.e. }G^d=1),
\qquad
\textbf{(C)}\ \ G=\Z/p^r,\ \ r<d.
\end{equation}
None of (B), (C) appears in \texttt{a\_wrong\_path.tex}; supplying them is the
whole game. \Cref{sec:object} fixes the object to be measured.
\Cref{sec:obs} records three structural observations (in particular: the
counterexample lives at $r=d$, and $r<d$ is essentially \emph{automatic} in the
totally ramified case). \Cref{sec:strategies} gives four attack strategies,
ranked. \Cref{sec:plan} proposes the smallest decisive computation.

% #############################################################################
\section{The object to be measured}\label{sec:object}
% #############################################################################

We keep the notation of \texttt{a\_wrong\_path.tex}: $G=\Z/p^r$ acting totally
ramified at $p_\infty$ (after the standard reductions), $H=\ker(G\to G/H)$ the
unique subgroup of order $p$ (the top Artin--Schreier layer), $G/H\cong
\Z/p^{r-1}$ the bottom layer, $f:C'\to C$ the $G/H$-cover, $p'_\infty\mapsto
p_\infty$. Blowing up $d\cdot p_\infty$ and $d\cdot p'_\infty$ gives the DVRs
$V_C\subset K(C^{(d)})$ and $V_{C'}\subset K(C'^{(d)})$ at the generic points
$\eta_C,\eta_{C'}$ of the exceptional divisors, and
\[
K_C=\Frac V_C,\quad
F=K\!\big((\cP_{G/H})^{[d]}\big),\quad
K_{C'}=\Frac V_{C'},\quad
E=K(\cP^{[d]_G}),\quad V_F=V_{C'}|_F .
\]

\begin{definition}[the $q$-defect]\label{def:defect}
Write $w:=V_{C'}$, $w_F:=V_F=w|_F$. The \textbf{ramification of $q$ along
$V_F$} is the ramification of the local extension
$\widehat{K_{C'}}/\widehat F$ at $w\mid w_F$: its inertia group
$I(w\mid w_F)$, its ramification index $e(w\mid w_F)$, its different
$\mathfrak d_{K_{C'}/F}(w\mid w_F)$, and (when wild and cyclic) its break
$\beta_q$ in upper numbering. We call any of these data the \emph{$q$-defect}.
\end{definition}

By \cref{eq:inj} and the diagram of \texttt{a\_wrong\_path.tex}, the Goal is
equivalent to surjectivity of \cref{eq:inj}, and the obstruction is precisely a
nontrivial $q$-defect that is ``aligned'' with $E/F$. Three reformulations of the
Goal, each a possible target:

\begin{enumerate}
\item \textbf{(Vanishing.)} $q$ is unramified along $V_F$, i.e.\
      $I(w\mid w_F)=1$. Then \cref{eq:inj} is an isomorphism and (H2)
      $\Rightarrow$ Goal immediately. \emph{Strongest; may be false in general
      --- see \Cref{sec:obs}\,(O2).}
\item \textbf{(Strict break separation.)} $E/F$ and $K_{C'}/F$ are cyclic
      wildly ramified of breaks $\beta_E,\beta_q$ at $V_F$ with
      $\beta_q\neq\beta_E$; more usefully $\beta_q<\beta_E$. Then
      $\widehat{K_{C'}}\neq\widehat E$, the local compositum
      $\widehat E\widehat{K_{C'}}$ is genuinely degree $p$ over
      $\widehat{K_{C'}}$ and ramified, contradicting (H2) unless $E/F$ was
      unramified to begin with. \emph{This is the cleanest sufficient
      condition; it directly kills the coincidence $\widehat{K_{C'}}=\widehat E$
      of \texttt{a\_wrong\_path.tex}.}
\item \textbf{(Conductor/different bound.)} $\mathfrak d_{K_{C'}/F}(w\mid w_F)$
      is too small to absorb $\mathfrak d_{E/F}(v_E\mid V_F)$ in the tower
      $F\subset K_{C'}\subset E\cdot K_{C'}$ versus $F\subset E\subset
      E\cdot K_{C'}$. \emph{Quantitative form of (2); good for an inequality
      proof.}
\end{enumerate}

\noindent The whole point is that all three are statements about $q$ \emph{alone}
(the $G/H$-cover and $d$), to be bounded by (B)+(C). They do not require
re-opening the absorption; they forbid it.

% #############################################################################
\section{Three structural observations}\label{sec:obs}
% #############################################################################

\paragraph{(O1) The counterexample sits exactly at the boundary $r=d$ --- (C)
excludes it.}
The fatal coincidence in \texttt{a\_wrong\_path.tex},
\Cref{ex:p2-counterexample}, is built with $G=\Z/p^2$ (so $r=2$) and the
degree-$2$ symmetric power ($d=2$): it is an $r=d$ phenomenon. The tame
companion in \texttt{inertia\_symmetry\_obstruction.tex} (eq.~``$d\cdot G_0$'')
fails the count $|G_0|\mid d$ at $d=|G_0|$, again the boundary. The hypothesis
$r<d$ in \cref{eq:hyp} is precisely the first regime \emph{not} covered by
either counterexample. So the right reading of the task is: \emph{prove the
obstruction is a boundary ($r\ge d$) phenomenon and disappears for $r<d$.}
This is encouraging --- (C) is not an arbitrary convenience; it is the
complement of the known bad locus.

\paragraph{(O2) Do not aim for ``$q$ unramified'' --- aim for a break
\emph{bound}.}
The leading-order (inertia / tame-model) computation of
\texttt{inertia\_symmetry\_obstruction.tex} already shows
$q$ is genuinely ramified along $V_{C'}$ in general: the diagonal of the
$E^d$-action survives, $I_{V_{C'}}(K_{C'}/F)\cong\Z/\gcd(d,|G/H|)$ in the tame
toy model, and $J\supseteq G^d$ in the wild model. So target~(1) of
\Cref{sec:object} is too much to hope for. The realistic targets are (2) and
(3): bound the \emph{break} of the $q$-defect strictly below that of $E/F$.
Higher ramification, not inertia, is the right invariant --- exactly the moral
of \texttt{inertia\_symmetry\_obstruction.tex}.

\paragraph{(O3) $r<d$ is essentially \emph{free} in the totally ramified case.}
For a totally ramified $\Z/p^r$ extension in which every layer is genuinely wild,
the upper-numbering breaks $u_1\le\cdots\le u_r$ are integers (Hasse--Arf) and
strictly increase, $u_1<u_2<\cdots<u_r$, whence $u_r\ge u_1+(r-1)\ge r$. The
bound (B) is $G^d=1$, i.e.\ $u_r<d$. Therefore
\[
r\;\le\;u_r\;<\;d ,
\]
so $r<d$ holds \emph{automatically} once $\cP$ is totally ramified with bound
$d$. Two consequences: (i) carrying (C) as a hypothesis through the induction
costs nothing in the case that matters; (ii) the genuinely dangerous regime
$r\ge d$ can only occur with \emph{coinciding} breaks $u_i=u_{i+1}$, i.e.\ when
some intermediate layer is \emph{not} wildly ramified --- which is the
degenerate situation the counterexamples exploit. This sharpens the goal: the
proof should \emph{use} the strict increase of the $u_i$.

% #############################################################################
\section{Four strategies}\label{sec:strategies}
% #############################################################################

\subsection*{Strategy A (recommended): compute the $q$-defect in the explicit
weighted-blowup local model, in terms of the breaks $u_i$ and $d$}

This is the direct continuation of the method that \emph{already works} for
$r=1$ (\Cref{cor:ramification_after_blowup_is_tame}, via
\Cref{theorem:artin_schreier_ramification} and the weighted blowup), and of the
local model set up in \texttt{inertia\_symmetry\_obstruction.tex}.

\smallskip
\noindent\emph{Set-up.} Take $k=\bar k$, $\widehat{K_C}=k((t))$, uniformiser $t$
at $p_\infty$. Present $\cP$ at $p_\infty$ as an Artin--Schreier--Witt
$\Z/p^r$-extension by a Witt vector $\vec a=(a_0,\dots,a_{r-1})$ over $k((t))$,
normalised so that the upper breaks are $u_1<\cdots<u_r<d$ (O3). The bottom
layer $G/H=\Z/p^{r-1}$ is given by $(a_0,\dots,a_{r-2})$ (breaks
$u_1,\dots,u_{r-1}$); the top layer $H=\Z/p$ adds $a_{r-1}$ (break $u_r$).
On the curve $C'$, a uniformiser $u$ at $p'_\infty$ satisfies
$t = u^{p^{r-1}}\cdot(\text{unit})$ (totally ramified $G/H$-cover).

\smallskip
\noindent\emph{Symmetric coordinates and the blowup valuation.} Use the $d$
copies $t_1,\dots,t_d$ (resp.\ $u_1,\dots,u_d$) and the elementary symmetric
functions $e_1,\dots,e_d$. By the thesis's own blowup computation
(\texttt{RamificationAfterBlowupIsTame.tex}, affine case:
$\nu_E(f)=\mathrm{ord}_0 f$) and the weighting of
\texttt{inertia\_symmetry\_obstruction.tex},
\[
\nu(e_i)=i ,\qquad \widetilde\nu(u_j)=1 ,
\]
and $V_C$ (resp.\ $V_{C'}$) is read off this weighted blowup. A pole $t^{-m}$
of a defining datum, rewritten in the $e_i$ and divided by the blowup scale,
acquires \emph{effective order} $\approx m/d$: this is exactly why $r=1$ works
($m<d\Rightarrow$ effective order $<1\Rightarrow$ the datum is a unit in the
residue field $\Rightarrow$ unramified at $\eta_C$).

\smallskip
\noindent\emph{What to compute.} The two completions at $w=V_{C'}$:
\begin{itemize}
  \item $\widehat E/\widehat F$: the symmetrised \emph{top} layer $a_{r-1}$.
        Extract its break $\beta_E$ at $V_F$ as a function of $u_r$ and $d$.
  \item $\widehat{K_{C'}}/\widehat F$ (the $q$-defect): produced by comparing
        the scheme symmetric power $(\cP_{G/H})^{(d)}$ with the torsor symmetric
        power $(\cP_{G/H})^{[d]}$ at $d\cdot p'_\infty$. Extract $\beta_q$ as a
        function of the \emph{bottom} breaks $u_1,\dots,u_{r-1}$ and $d$.
\end{itemize}

\noindent\emph{Target inequality.} Show
\begin{equation}\label{eq:target}
\boxed{\ \beta_q \;<\; \beta_E \quad\text{whenever}\quad u_r<d\ \text{and}\ r<d.\ }
\end{equation}
By target~(2) of \Cref{sec:object}, \cref{eq:target} forbids
$\widehat{K_{C'}}=\widehat E$ and makes (H2) imply the Goal. The heuristic in
its favour: $\beta_E$ is driven by the \emph{largest} break $u_r$, while
$\beta_q$ is driven only by the \emph{bottom} breaks $u_{\le r-1}$ (the map $q$
sees only $\cP_{G/H}$), and by (O3) $u_{r-1}<u_r$ strictly. The role of $r<d$ is
to keep the Witt-vector \emph{carries} (which can lift effective orders when the
$d$ branches collide) below the blowup scale $1$, so that no accumulation of the
$r$ bottom layers reaches $\beta_E$.

\smallskip
\noindent\emph{Why this is the recommended route.} (i) It is the proven method
for $r=1$, pushed one Witt-degree further. (ii) It computes
$I(v_E\mid V_C)$ \emph{directly} and never base-changes along $q$, so the
absorption mechanism cannot even be expressed --- it is structurally immune to
the \texttt{a\_wrong\_path} obstruction. (iii) It produces a number ($\beta_q$),
which is what every other strategy ultimately needs.

\smallskip
\noindent\emph{Risk.} The Witt-vector carry analysis in $d$ symmetric variables
is the technical heart and is genuinely involved; the clean statement
\cref{eq:target} must be checked, not assumed. \Cref{sec:plan} isolates the
smallest instance.

\subsection*{Strategy B: conductor / different accounting in the parallelogram}

Stay abstract and use only the transitivity of the different. In the two towers
\[
\widehat F\subset\widehat E\subset \widehat E\widehat{K_{C'}}
\qquad\text{and}\qquad
\widehat F\subset\widehat{K_{C'}}\subset \widehat E\widehat{K_{C'}},
\]
transitivity gives
$\mathfrak d_{E K_{C'}/F}
   =\mathfrak d_{E K_{C'}/E}+\mathfrak d_{E/F}
   =\mathfrak d_{E K_{C'}/K_{C'}}+\mathfrak d_{K_{C'}/F}$
(pulled back appropriately). Hypothesis (H2) is
$\mathfrak d_{E K_{C'}/K_{C'}}=0$. Hence
\begin{equation}\label{eq:diff}
\mathfrak d_{E/F}\ \big(\text{pulled up}\big)
   \;=\;\mathfrak d_{E K_{C'}/E}\;-\;\mathfrak d_{K_{C'}/F}\ \big(\text{pulled up}\big).
\end{equation}
The Goal is $\mathfrak d_{E/F}=0$. So one must show the right-hand side of
\cref{eq:diff} vanishes, i.e.\ that base-changing $K_{C'}/F$ up to $E$ does not
\emph{increase} its different: $\mathfrak d_{E K_{C'}/E}=\mathfrak
d_{K_{C'}/F}\cdot(\text{pullback})$. This is automatic when $E/F$ and
$K_{C'}/F$ are arithmetically disjoint, e.g.\ when their breaks differ ---
recovering \cref{eq:target} as the clean sufficient condition, now phrased via
the different. Combine with the Hasse--Arf/Herbrand convexity already sketched
in \texttt{inertia\_symmetry\_obstruction.tex}\,\S(higher):
\[
\varphi_G(m_r)=\varphi_{G/H}(m_{r-s})+\tfrac{\varphi_H(m_r)-m_{r-s}}{p^{\,r-s}}
\quad\text{(convex combination)}\ \Rightarrow\ \varphi_G(m_r)<d ,
\]
to bound the \emph{total} conductor of $\cP$ by $d$ and split it across the
layers. \emph{Pro:} no symmetric-variable bookkeeping; rests on standard
ramification calculus. \emph{Con:} the disjointness/non-increase step is exactly
where the absorption hides --- you still need an input equivalent to
\cref{eq:target}, so B is best as the \emph{packaging} of A's output, not a
replacement for it.

\subsection*{Strategy C: rule out $\widehat{K_{C'}}=\widehat E$ by a global
break invariant}

Attack target~(2) head-on without computing $\beta_q$ exactly: find a
\emph{global} invariant of $K_{C'}/F$ at $V_F$ that is forced by (B)+(C) to differ
from the corresponding invariant of $E/F$. Candidates:
\begin{itemize}
  \item \textbf{Source of the break.} $\beta_E$ is a break of the $H$-layer
        (top), $\beta_q$ a break of the symmetrised $G/H$-tower (bottom). If one
        shows $\beta_q\in\{$values produced by $u_1,\dots,u_{r-1}\}$ and
        $\beta_E\in\{$values produced by $u_r\}$, then (O3)'s strict increase
        $u_{r-1}<u_r$ separates them.
  \item \textbf{Divisibility / residue.} Even if breaks coincide numerically,
        the leading coefficients (the ``$u^{-m}$ change of variable'' of
        \Cref{theorem:artin_schreier_ramification}) may be forced distinct by
        the explicit form of $q$. A coincidence $\widehat{K_{C'}}=\widehat E$
        requires equality of Artin--Schreier classes in
        $\widehat F/\wp(\widehat F)$; show $(B)+(C)$ make the class of
        $K_{C'}/F$ land in a different coset than that of $E/F$.
\end{itemize}
\emph{Pro:} potentially a short, conceptual proof. \emph{Con:} ``find the right
invariant'' is open-ended; likely needs A's local model to even identify the
invariant.

\subsection*{Strategy D (fallback): change the resolution to make $\eta_{C'}$
codimension one for $q$}

The reason the existing tool fails is purely that
\Cref{lemma:torsor_unramified_codim_one} controls $q$ only in
\emph{codimension one} on $C'^{(d)}$, whereas $V_{C'}$ sits over the
codimension-$d$ point $d\cdot p'_\infty$. Idea: choose a modified
blowup / weighted resolution $X_{C'}'\to C'^{(d)}$ in which the chosen valuation
$\eta_{C'}$ \emph{is} a codimension-one point already on the source of $q$ (or
factor $q$ through a sequence of codimension-one steps), then apply the
codimension-one lemma iteratively. This trades a hard local computation for a
careful birational geometry of the symmetric power along the small diagonal.
\emph{Pro:} reuses a lemma you already have. \emph{Con:} you must check the
modified resolution still computes the \emph{same} $V_C\hookrightarrow V_{C'}$
and still detects the modulus bound $d$; the weighted blowup was chosen for a
reason (the unweighted one ``does not even extend $V_C$ to $V_{C'}$'', per
\texttt{inertia\_symmetry\_obstruction.tex}).

% #############################################################################
\section{Recommended next step: the decisive small computation}\label{sec:plan}
% #############################################################################

Everything hinges on one number, $\beta_q$. Compute it in the first open case
and check \cref{eq:target}:
\begin{quote}
\textbf{Test case.} $k=\bar{\F}_p$, $G=\Z/p^2$ ($r=2$), $H=\Z/p$ bottom,
$G/H=\Z/p$. Take $\cP$ totally ramified at $p_\infty$ with the bottom-layer
break $u_1$ and top-layer break $u_2$, $u_1<u_2$. Choose $d$ with
$u_2<d$ and $r=2<d$ --- the smallest honest choice is $d=3$ (so the
\texttt{a\_wrong\_path} counterexample's $d=2=r$ is just excluded). Then:
\begin{enumerate}
  \item Write the $G/H=\Z/p$ cover $C'\to C$ locally as $y^p-y=t^{-u_1}$;
        $u=$ uniformiser at $p'_\infty$, $t=u^{p}\cdot\text{unit}$.
  \item Form $C'^{(3)}$, the elementary symmetric functions $e_1,e_2,e_3$ of the
        three branches, and the weighted blowup $\nu(e_i)=i$. Identify $V_{C'}$,
        $V_F=V_{C'}|_F$, and the local extension $\widehat{K_{C'}}/\widehat F$
        induced by $q$. \textbf{Read off $\beta_q$.}
  \item Symmetrise the top layer $a_1$ (break $u_2$) the same way; read off
        $\beta_E$.
  \item Check $\beta_q<\beta_E$. If it holds, (H2) $\Rightarrow$ Goal for
        $(r,d)=(2,3)$; if it fails, the failure pattern tells you whether the
        true hypothesis is $r<d$, or $p^{\,?}\mid d$, or something finer.
\end{enumerate}
\end{quote}

\noindent A single clean computation here either (a) proves the first case and
gives the template for the induction, or (b) produces a new, $r<d$
counterexample that pins down the genuinely required hypothesis --- both are
decisive outcomes. It is the cheapest experiment with the highest information
content, and it directly resolves the red TODO at
\texttt{RamificationAfterBlowupIsTame.tex:671}.

\paragraph{Summary of the logic.}
\begin{enumerate}
  \item Goal $\equiv$ surjectivity of \cref{eq:inj} $\equiv$ no $q$-defect
        aligned with $E/F$.
  \item The missing input is a bound on the $q$-defect by $(\cP,d)$; (B) and
        (C) are the unused hypotheses that can give it.
  \item The known obstructions are $r=d$ boundary phenomena (O1), and $r<d$ is
        free in the case that matters (O3); so the obstruction should vanish.
  \item The decisive quantity is the break $\beta_q$; compute it (Strategy A,
        \Cref{sec:plan}), package the conclusion via the different (Strategy B),
        and the absorption is forbidden.
\end{enumerate}

\printbibliography

\end{document}
